\chapter{The second track}
\label{ch:talent}

\section{The claim, stated carefully}

This is the most speculative chapter in the book and it should announce that
before it argues anything.

The claim is that Cuba can build a second export sector, alongside tourism, out
of educated people selling services across a border without leaving the island
--- software, engineering, design, back-office work, data work, Spanish-language
content, remote professional services --- and that this sector could plausibly
reach the same order of magnitude as tourism within a decade.

The claim is not that Cuba becomes an artificial-intelligence power. It will not.
Training large models requires electricity Cuba does not have, capital it cannot
raise, and hardware it cannot legally buy. Any proposal that has Cuba building
data centres is not a proposal about Cuba. What Cuba can plausibly do is supply
\emph{labour} into the global services market that the current technology wave
has expanded --- the annotation, evaluation, localisation, integration, tooling
and ordinary software engineering that sit around the models rather than inside
them --- and it can do so in Spanish, in the New York time zone, with a
workforce that can read.

The arithmetic is worth putting on the page because it is the reason the chapter
exists. A hundred thousand Cubans --- about two per cent of the labour force ---
earning an average of eighteen thousand dollars a year in remote services would
generate roughly \$1.8 billion annually in foreign exchange.\unv{} That is the
same order as Cuban tourism earnings in a good year, from a sector requiring no
hotels, no beach, no water supply, no fuel and almost no capital. The per-worker
figure is well below what a comparable worker earns in a high-income country and
well above anything available in Cuba today, which is precisely why the market
would clear.

\section{Why it has not happened}

Cuba has had educated, underemployed, English-capable technical workers for
thirty years, and the global market for remote services has existed for twenty.
The sector did not appear. The reasons are specific, and identifying the binding
one is the whole content of this chapter.

It is not talent. Cuban technical education is the asset Chapter~\ref{ch:socialstate}
described, and Cuban programmers working abroad are competitive.

It is not, primarily, the embargo. American clients are constrained, but Spain,
Mexico, Argentina, Canada, Germany and the entire Spanish-speaking world are
not, and none of them buys Cuban services at scale either.

It is not, mainly, connectivity --- though connectivity is a real constraint and
the next section is about it. Plenty of countries export services on worse
bandwidth than Havana has had since 2018.

\emph{The binding constraint is that a Cuban cannot lawfully and reliably
contract with a foreign client and be paid into an account they control.} Until
2021 a Cuban private operator had no legal personality and so could not sign a
contract as an entity. The permitted-activity list was written in the vocabulary
of trades and had no category for a software developer or a design studio.
Foreign-currency accounts were restricted. International card and transfer rails
did not reach individuals. And professionals trained at state expense were, for
much of the period, prohibited from practising their profession privately at all.

The result is that Cubans who do this work do it informally, through a relative's
foreign account, or through an intermediary in Miami or Madrid who takes a cut
and holds the client relationship --- or, increasingly, they emigrate, which
converts a potential export sector into a permanent loss.

This is the good news in the chapter. The binding constraint is legal, it is
domestic, it costs the state nothing to remove, and it does not require anyone's
permission.

\begin{design}{The right to sell your own work}
\label{d:contract}
\begin{rules}
\rl{1}{The core right} Any resident may contract directly with a foreign client,
in any currency, without licence, intermediary or prior authorisation, and may
receive payment into an account in their own name and spend from it. This is the
single most important clause in Part~III's growth strategy and it is a repeal,
not a programme.

\rl{2}{Professions included} Explicitly extended to professions trained at state
expense --- engineers, physicians, architects, lawyers, scientists. The exclusion
of professionals from self-employment since 1993 is the reason Cuba's private
sector is a collection of trades, and it is the reason its most valuable workers
had only two options, the state or the airport.

\rl{3}{No enumerated list} Per Design~\ref{d:commit}.2, service exports are
permitted by right. Any list of permitted activities in this domain would be
obsolete before it was published, which is a general argument against lists and a
decisive one here.

\rl{4}{A simple tax that is worth paying} Income from cross-border services pays
the low flat rate of Design~\ref{d:tax}.6, filed annually, with no deductions and
no separate social contribution. A mobile base taxed heavily is not taxed; a
mobile base taxed lightly and simply is registered, counted, and eventually
becomes a constituency for the system that registered it.

\rl{5}{Equipment} Computers, phones, networking hardware and their parts enter
free of duty and licence, per Design~\ref{d:genrights}.6. Cuba currently has
citizens buying their own tools with remittance money and a customs regime that
treats it as smuggling.

\rl{6}{Classification} Cross-border services are classified and counted as
exports in the national accounts and the balance of payments, which sounds
clerical and is not: what is not measured is not defended in the budget, and this
sector will need defending against the ministries whose staff it draws.
\end{rules}
\end{design}

\section{Connectivity}

Cuba came online late and expensively. Mobile data arrived in December 2018;
international capacity has run principally through a single submarine cable laid
from Venezuela in 2011, which is both a bandwidth constraint and a single point
of failure; the state telecommunications company is a monopoly; and the price of
data as a share of Cuban income has been among the highest in the world, which is
what the protests over the 2025 tariff changes were about.

No second track exists on that basis. A foreign client contracting with a Cuban
supplier is buying reliability, and reliability is exactly what a single cable, a
monopoly operator, a fragile grid and a state that switched off mobile data
nationally on 11 July 2021 do not provide.

\begin{design}{Connectivity as infrastructure}
\label{d:connect}
\begin{rules}
\rl{1}{Redundancy} At least two additional international routes, on diverse
physical paths and diverse landing points, procured on published competitive
terms. One cable is not a network.

\rl{2}{Open wholesale access} Whatever the ownership of the international and
backbone capacity, access to it is sold to any licensed retail provider on
published, non-discriminatory terms, regulated by an independent authority on
the model of Design~\ref{d:genrights}.4. Retail competition is what moves the
price; full privatisation of the incumbent is not required and should not be
attempted early.

\rl{3}{No political interruption} Per Design~\ref{d:minimum}.5. This is
restated here because it is a commercial clause as much as a civil-liberties
one: a country whose government has demonstrated that it will switch off the
internet during a political event cannot sell service-level agreements, and
every potential client knows it.

\rl{4}{Price transparency} Published wholesale and retail tariffs, and a
published affordability metric --- data cost as a share of median income ---
reported annually.

\rl{5}{Power} Telecommunications sites are critical loads under
Design~\ref{d:energyseq}.1 and are among the first to receive photovoltaics and
storage. A network that goes down with the grid is not a network either.
\end{rules}
\end{design}

\section{The residence regime}

Only after the two designs above does a visa policy make any sense, and the
ordering matters because the international discussion of this subject has it
backwards.

The digital-nomad visa is a modest instrument. The honest reading of the
international experience --- Estonia, Portugal, Barbados, Croatia, Costa Rica ---
is that these schemes attract a few thousand people, generate useful but small
consumption spending, and do not by themselves build an industry. Barbados's
Welcome Stamp was a competent response to a pandemic-era collapse in tourism, not
an economic transformation. Estonia's e-Residency, the most cited of them, has
enrolled a large number of people and produced a much smaller number of real
businesses paying real Estonian tax.

What actually built service exports in the comparator countries was not visas for
foreigners. It was \emph{residents being allowed and able to sell their work
abroad}, plus predictable law and a payment system --- which is
Design~\ref{d:contract}. Uruguay's software export sector, the most relevant Latin
American case, was built on domestic firms and domestic engineers, not on
inbound nomads.

So the residence regime here is proposed for three narrower and more defensible
reasons: it brings in people who train and hire locally; it brings in demand for
the long-stay accommodation that Chapter~\ref{ch:tourism} wants to grow instead
of hotel rooms; and --- per Design~\ref{d:pension}.3 --- more working-age
residents is the cheapest available instrument for a dependency ratio.

\begin{design}{Residence for mobile workers}
\label{d:residence}
\begin{rules}
\rl{1}{The permit} A renewable residence permit for a person with foreign income
from remote work or a foreign business, granted on published objective criteria
--- income threshold, health cover, clean record --- with a statutory decision
period and deemed grant if the period passes. No discretion, no interview, no
sponsor.

\rl{2}{Tax certainty} A published flat rate on locally-remitted income for a
fixed number of years, legislated rather than administered, with the treatment
of foreign income stated in advance. Uncertainty deters this population far more
than rates do.

\rl{3}{Ordinary rights} The right to rent, to open an account, to import personal
equipment, to bring family, and to use the health system on the same terms as
residents while paying the same taxes. A permit that does not allow a bank
account is not a permit.

\rl{4}{A route to staying} A defined path to permanent residence and, for those
who want it, naturalisation. The objective is residents, not visitors; a scheme
that resets every year selects for people who do not put down roots.

\rl{5}{Local linkage} A reduced rate for permit-holders who employ or formally
train Cuban residents, verified simply and audited. This is what converts an
inbound flow into a domestic industry, and it is the difference between this
proposal and a beach with better wifi.
\end{rules}
\end{design}

\section{The diaspora, and not taxing the door}

Two million people abroad, most of them one flight away, many with capital and
professional standing, and a substantial number who would work with Cuba if
working with Cuba were possible. They are simultaneously the natural first
clients for this sector, the natural first investors under
Chapter~\ref{ch:capital}, and the largest single reservoir of the working-age
people Chapter~\ref{ch:premises} says the country has lost.

\begin{design}{Circulation, not retention}
\label{d:diaspora}
\begin{rules}
\rl{1}{No exit tax} Cuba does not tax, fine, or impose a training levy on
emigrants. The instrument is popular with governments losing skilled workers, it
collects almost nothing, it is trivially avoided by simply not returning, and it
converts a temporary departure into a permanent one. It is also, in Cuba
specifically, a tax on the grievance this book is trying to settle.

\rl{2}{Return without penalty} Full restoration of residence, property rights,
professional registration and pension entitlement for returning Cubans,
automatic and without discretionary review, including for those who left
irregularly or abandoned a mission.

\rl{3}{Dual nationality} Recognised, with the right to hold property, to invest,
to work, to vote in local elections after a residence qualification, and to
inherit. Chapter~\ref{ch:capital} depends on this; so does the pension
portability of Design~\ref{d:pension}.4.

\rl{4}{Voluntary contribution} Cubans abroad may contribute voluntarily to the
Cuban social insurance system and accrue entitlement, per
Design~\ref{d:pension}.4. This converts a diaspora into a contributor base,
which no other instrument does.

\rl{5}{Faculty and clients} Programmes to bring diaspora academics and
professionals back on short teaching, clinical and advisory terms, with no
political test, administered by universities rather than by ministries.
\end{rules}
\end{design}

\section{The universities}

The sector proposed here needs a pipeline, and Chapter~\ref{ch:socialstate}
identified the specific gap: Cuba produces excellent engineers, physicians and
mathematicians and very few people trained in the disciplines an institutional
transition and a services economy actually run on --- commercial law, accounting,
finance, management, applied economics, and the empirical social sciences.

\begin{design}{Academic capacity}
\label{d:university}
\begin{rules}
\rl{1}{Autonomy} Universities appoint their own staff, set their own curricula
and admit their own students on published academic criteria. Political record is
not an admission or appointment criterion, and the prohibition is justiciable.

\rl{2}{The missing disciplines} A funded programme to build law, accounting,
finance, management and applied economics faculties, with foreign partnership,
visiting appointments from the diaspora, and professional accreditation
recognised internationally. This is a decade-long constraint on everything in
Part~III --- Design~\ref{d:courts}.6 needs commercial judges, and
Chapter~\ref{ch:welfare} needs tax administrators --- and it binds from year one.

\rl{3}{Foreign partnership} Joint programmes, dual degrees and branch
arrangements with foreign universities, permitted by right and regulated for
quality only.

\rl{4}{Publish and be read} Cuban academic work published openly, and Cuban
academics free to publish abroad, collaborate and travel. The dismantling of the
Centre for the Study of the Americas in 1996 taught Cuban social science what
happens to research that analyses policy, and that lesson has to be publicly
unlearned before anyone will do the work this book's later phases need.
\end{rules}
\end{design}

\section{What this chapter cannot promise}

Three honest limits.

The sector is small until the electricity is fixed. Everything here is downstream
of Chapter~\ref{ch:energy}, and a remote worker in a province with sixteen-hour
blackouts cannot take a call.

The sector competes for exactly the people the state most needs. Every engineer
who moves to remote work at eighteen thousand dollars a year is an engineer not
in the ministry, the utility or the university at a state wage --- which is why
Design~\ref{d:pay} is a precondition for this chapter and not a separate topic.

And the sector is a rival to emigration, not a complement to it. Its real
justification is that it offers a person the income they left for without
requiring them to leave. If it works, the outflow slows; if the outflow does not
slow, it is not working, and Chapter~\ref{ch:sequence} makes net migration one of
the indicators by which this book can be judged to be failing.
